\chapter{Modellbewertung}\label{chap:modellbewertung}

\section{Versuchsserie A01--A09}

Um verschiedene Modellkonfigurationen zu vergleichen, wurde die aktuelle Versuchsserie A01 bis A09 durchgeführt. 
Dabei wurden unterschiedliche Strategien, Profilbreiten und Regularisierungsparameter getestet. Ziel davon ist es, die Einflüsse einzelner Parameter auf die Modellleistung zu analysieren, \parencite{LTR_Configuration_Test_Series}.

\subsection{Untersuchte Konfigurationsparameter}

Die verschiedenen Varianten der Konfiguration haben unterschiedliche Auswirkungen auf die Trainingslabels, die Kandidatenprofile sowie die Modellregularisierung. 
Die untersuchten Konfigurationsparameter werden in Tabelle \ref{tab:konfigurationsparameter} erläutert.

\begin{table}[H] 
    \centering 
    \caption{Untersuchte Konfigurationsparameter} 
    \label{tab:konfigurationsparameter} 
    \begin{tabularx}{\textwidth}{p{4cm}X} 
        \toprule \textbf{Parameter} & \textbf{Bedeutung} \\ 
        Same-Family-Gewicht &  Negative Kandidaten aus derselben Makrofamilie werden innerhalb der Verlustfunktion unterschiedlich gewichtet. Dadurch kann die Variantenauflösung stärker betont werden. \\ 
        Profilbreite & Begrenzt die Anzahl der gespeicherten Profilwerte innerhalb eines Kandidatenprofils. Kleinere Profilbreiten erzeugen eine kompaktere Kandidatenrepräsentation. \\ 
        L1-Regularisierung ($\alpha$) & Bestraft große Modellgewichte und reduziert die Modellkomplexität. Dadurch werden irrelevante Merkmale stärker unterdrückt. \\ 
        L2-Regularisierung ($\lambda$) & Stabilisiert die Gewichtung der einzelnen Merkmale und reduziert die Gefahr einer Überanpassung an einzelne Trainingsbeispiele. \\ 
        Min. Split Gain & Definiert den minimal notwendigen Informationsgewinn, um eine Baumaufteilung im LightGBM-Modell zu ermöglichen. Höhere Werte führen dadurch zu konservativeren Entscheidungsbäumen. \\ 
        \bottomrule 
    \end{tabularx} 
\end{table} 

Bei der L1-Regularisierung wird die Summe der Absolutwerte der Modellgewichte bestraft:
\begin{align}
  \lambda_1 \sum_j |w_j|.
\end{align}

Dadurch werden einzelne Gewichte auf null gesetzt. 
Die L2-Regularisierung bestraft dagegen die Summe der quadrierten Gewichte:
\begin{align}
  \lambda_2 \sum_j w_j^2.
\end{align}

Dadurch werden besonders große Gewichte begrenzt. Bei LightGBM beziehen sich diese Regularisierungsstrafen hauptsächlich auf die Ausgabewerte der Baumblätter.

A02 bis A07 wurden so gewählt, dass immer nur ein Parameter gezielt verändert wurde. Dadurch lassen sich die Auswirkungen der einzelnen Konfigurationsparameter isoliert untersuchen.
Bei A08 und A09 werden mehrere Parameter (L1 und L2) gleichzeitig verändert, da diese bei LightGBM konzeptionell denselben Hebel darstellen (Bestrafung der Blattgewichte). 
Um eine stärkere Regularisierung zu erzielen, ist es in der Praxis sinnvoll, beide gemeinsam hochzudrehen, nicht nur einen isoliert.
Bei A09 kommt zusätzlich ein positiver Wert für \texttt{min\_split\_gain} zum Einsatz, um dadurch einen noch größeren Einfluss auf eine stärkere Regularisierung zu erzielen, \parencite{LTR_Configuration_Test_Series}.

\subsection{Referenzkonfiguration A01}

A01 dient als Referenzkonfiguration und bildet die Basis für die Vergleichsexperimente. Sie enthält die Standardwerte für alle untersuchten Konfigurationsparameter.
Alle weiteren Konfigurationen A02 bis A09 werden gegenüber dieser bewertet.

Die Referenzparameter sind in Tabelle \ref{tab:a01_referenz} zusammengefasst.

\begin{table}[H] 
    \centering 
    \caption{Referenzkonfiguration A01} 
    \label{tab:a01_referenz} 
    \begin{tabular}{ll} 
        \toprule 
        \textbf{Parameter} & \textbf{Wert} \\ 
        \midrule 
        Exaktes Targettraining & Ja \\ 
        Same-Family-Gewicht & 5,0 \\ 
        Profilbreite & 16 \\ 
        L1-Regularisierung ($\alpha$) & 0,05 \\ 
        L2-Regularisierung ($\lambda$) & 0,2 \\ 
        Min. Split Gain & 0,0 \\ 
        \bottomrule 
    \end{tabular} 
\end{table}

Ausgehend von dieser Basis werden die Varianten der Konfigurationen A02 bis A09 verändert.
Tabelle \ref{tab:konfigurationsabweichungen} zeigt die gegenüber A01 vorgenommenen Veränderungen. 

\begin{table}[H] 
    \centering 
    \caption{Abweichungen der Konfigurationen gegenüber der Referenz A01} 
    \label{tab:konfigurationsabweichungen} 
    \begin{tabular}{lp{8cm}} 
        \toprule 
        \textbf{Konfiguration} & \textbf{Änderung gegenüber A01} \\ 
        \midrule 
        A02 & Same-Family-Gewicht von 5,0 auf 2,0 reduziert \\ 
        A03 & Same-Family-Gewicht von 5,0 auf 8,0 erhöht \\ 
        A04 & Same-Family-Gewicht von 5,0 auf 12,0 erhöht \\ 
        A05 & Hierarchische Relevanzlabels statt ausschließlich exakter Ziel-Makros \\ 
        A06 & Profilbreite von 16 auf 12 reduziert \\ 
        A07 & Profilbreite von 16 auf 8 reduziert \\ 
        A08 & L1-Regularisierung auf 0,1 und L2-Regularisierung auf 0,5 erhöht \\ 
        A09 & L1-Regularisierung auf 0,5, L2-Regularisierung auf 1,0 und \texttt{min\_split\_gain} auf 0,01 erhöht \\ 
        \bottomrule 
    \end{tabular} 
\end{table}

Dadurch lässt sich der Einfluss der jeweiligen Parameteränderung der entsprechenden Konfiguration zuordnen, \parencite{LTR_Configuration_Test_Series}.

\subsection{Ergebnisse}

Tabelle \ref{tab:konfigurationsvergleich} zeigt die getesteten Modellkonfigurationen und deren Leistung im Vergleich.

\begin{table}[H] 
    \centering 
    \caption{Ergebnisse der A01--A09-Konfigurationen auf dem Testdatensatz} 
    \label{tab:konfigurationsvergleich} 
    \small 
    \begin{tabular}{lcccc} 
        \toprule 
        \textbf{Konfiguration} & \textbf{Recall@1} & \textbf{Recall@3} & \textbf{Recall@10} & \textbf{MRR} \\ 
        \midrule 
        A01 Referenz & 0,3241 & 0,4822 & 0,6798 & 0,4418 \\ 
        A02 Gewicht 2 & 0,3162 & 0,4862 & 0,6680 & 0,4399 \\ 
        A03 Gewicht 8 & 0,3241 & 0,4743 & 0,6838 & 0,4397 \\ 
        A04 Gewicht 12 & 0,2925 & 0,4783 & 0,6759 & 0,4195 \\ 
        A05 Hierarchisch & 0,3162 & 0,4783 & 0,6719 & 0,4328 \\ 
        A06 Profilbreite 12 & 0,3162 & 0,4862 & 0,6877 & 0,4374 \\ 
        A07 Profilbreite 8 & 0,3241 & 0,4862 & 0,6719 & 0,4439 \\ 
        A08 Regularisierung mittel & \textbf{0,3478} & 0,4783 & 0,6798 & \textbf{0,4534} \\ 
        A09 Regularisierung stark & 0,3399 & 0,4783 & \textbf{0,6917} & 0,4505 \\ 
        \bottomrule 
    \end{tabular} 
\end{table}

\subsection{Same-Family-Gewicht}

Mit den Konfigurationen A02 bis A04 wird der Einfluss des Same-Family-Gewichts untersucht. 
Die Ergebnisse zeigen, dass eine Verringerung des Gewichts auf 2 den Recall@1 auf 31,62\,\% senkt, während eine Erhöhung auf 8 keine Auswirkung zeigt.
Die Gewichtung auf 12 führt zu einer deutlichen Verschlechterung des Recall@1 auf 29,25\,\%.

Daraus lässt sich schließen, dass ein zu hohes Same-Family-Gewicht die Lernstabilität negativ beeinflusst.

\subsection{Hierarchische Labeldefinition}

A05 untersucht die hierarchische Relevanzdefinition. Makros \parencite{EPLAN-Makros} derselben Familie und desselben Präfixes werden dabei hierarchisch relevant betrachtet.
Es zeigt sich jedoch keine Verbesserung gegenüber dem exakten Targettraining. Der Recall@1 sinkt von 32,41 \,\% auf 31,62 \,\% und die MRR von 44,18 \,\% auf 43,28 \,\%.
Die dadurch entstandene reichhaltigere Zieldefinition hat somit keinen positiven Einfluss auf die Modellleistung, \parencite{LTR_Configuration_Test_Series}.

\subsection{Profilbreite}

In den Konfigurationen A06 und A07 werden unterschiedliche Profilbreiten untersucht. Bei A06 beträgt diese zwölf, bei A07 acht.
Die Erhöhung der Profilbreite führt zu einer Verschlechterung des Recall@1 von 32,41 \,\% auf 31,62 \,\% und der MRR von 44,18 \,\% auf 43,74 \,\%.
Eine Verringerung der Profilbreite führt zu keiner nennenswerten Veränderung der Modellleistung. 
Die Ergebnisse zeigen, dass eine verringerte Profilbreite die Modellleistung nicht signifikant verbessert, \parencite{LTR_Configuration_Test_Series}.

\subsection{Regularisierung}

Die deutlichsten Verbesserungen werden durch die Veränderung der Regularisierung erzielt.
A08 verwendet eine mittlere L1- und L2-Regularisierung (\(\alpha = 0{,}1\), \(\lambda = 0{,}5\)). 
Dadurch wird mit einem Recall@1 von 0,3478 der höchste Top-1-Wert aller untersuchten Konfigurationen erreicht. Auch der höchste MRR-Wert von 0,4534 wird durch A08 erreicht.
Mit A09 wird die Regularisierung nochmals deutlich erhöht (\(\alpha = 0{,}5\), \(\lambda = 1{,}0\)). Zusätzlich wird ein positiver Wert für \texttt{min\_split\_gain} gesetzt, um die Stabilität der Baumaufteilung zu erhöhen. 
Diese Konfiguration erreicht den höchsten Recall@10 von 0,6917 und die zweithöchste MRR von 0,4505.

Die L1/L2-Regularisierung wirkt sich bei LightGBM direkt auf die Blattgewichte aus.
Sie dämpft diese bei geringer Aussagekraft auf null und unterdrückt dadurch das Rauschen der 2.381 Merkmale bei nur 984 Trainingsbeispielen, ohne das schwache Unterscheidungssignal zwischen Makrofamilien-Varianten zu verlieren.
Bei A09 wird durch die stärkere Dämpfung auch dieses Signal mit reduziert, was die Top-1-Präzision verringert, dabei aber die Gesamtrangfolge robuster macht, \parencite{LTR_Configuration_Test_Series}.

\subsection{Top-Konfiguration}

Die beste Konfiguration der Studie ist A08, da sie sowohl den höchsten Recall@1 als auch den höchsten MRR erreicht\parencite{LTR_Configuration_Test_Series}.

Die wichtigsten Kennzahlen dieser Variante sind in Tabelle \ref{tab:a08_best} dargestellt. 

\begin{table}[H] 
    \centering 
    \caption{Ergebnisse der besten Konfiguration A08} 
    \label{tab:a08_best} 
    \begin{tabular}{lr} 
        \toprule 
        \textbf{Kennzahl} & \textbf{Wert} \\ 
        \midrule Recall@1 & 0,3478 \\ 
        Family Recall@1 & 0,8735 \\ 
        Recall@10 & 0,6798 \\ 
        Family Recall@10 & 0,9368 \\
        MRR & 0,4534 \\ 
        Mittlerer Zielrang & 23,96 \\ 
        90\,\%-Perzentil & 61,6 \\ 
        Schlechtester Rang & 360 \\ 
        \bottomrule 
    \end{tabular} 
\end{table}

Die Fehleranalyse dieser Konfiguration zeigt, dass 7,9\,\% aller Anfragen einer fachlich falschen Makrofamilie zugeordnet wurden.
52,6\,\% der Fehlklassifikationen erfolgen innerhalb der korrekten Makrofamilie. 
Die größte Herausforderung liegt damit weiterhin in der Variantenauflösung technisch ähnlicher Makrofamilien, \parencite{LTR_Configuration_Test_Series}.

